\documentclass[10pt,a4paper]{amsart}
\usepackage[margin=19mm]{geometry}
\usepackage{amsmath,amssymb,mathtools}
\usepackage[scr=rsfs,cal=euler]{mathalfa}
\usepackage{xfrac}
\usepackage[unicode,hidelinks,bookmarks=false]{hyperref}
\newcommand{\ie}{i.e.\ }
\newcommand{\cf}{cf.\ }
\newcommand{\resp}{respectively\ }
\newcommand{\Subsection}{\subsection}
\providecommand{\nobreakdash}{\nobreak\hbox{-}}
\setlength{\emergencystretch}{2em}
\multlinegap0pt
\usepackage{amssymb}
\usepackage{tikz}
\usepackage[normalem]{ulem}
\usepackage{enumerate}
\usepackage{mathtools}
\newtheorem{claim}{Claim}
\title{The unbreakable quasi-graphic matroids}
\author{Sayantani Bhattacharya, John David Clifton, Zach Walsh}
\date{Source: arXiv:2605.12811v1 (CC BY 4.0)}
\begin{document}
\maketitle

\noindent\emph{Source and licence.} The unbreakable quasi-graphic matroids. Version arXiv:2605.12811v1 (CC BY 4.0), distributed by arXiv under the Creative Commons Attribution 4.0 International licence, http://creativecommons.org/licenses/by/4.0/, as recorded in the arXiv licence metadata for this exact version and shown on the abstract page https://arxiv.org/abs/2605.12811v1. Full licence notice: \texttt{licenses/arxiv-2605.12811-CC-BY-4.0.txt}.\par\smallskip

\noindent\textit{Editorial scope.} Counted unit: the article's claim claim (source0.tex lines 1358--1360). The article's complete original proof of it is delivered with it (source0.tex lines 1361--1382, one \texttt{proof} environment of 22 lines). No mathematics is written, repaired or re-derived: the statement and the proof are the article's own lines, in the article's own notation, and no retained line is substituted or re-flowed.  No further source-local context is needed: every cross-reference of every retained line resolves inside the two delivered spans.  Nothing was omitted from any retained span, so the package carries no omission record.  No retained line contains a citation, so the wrapper carries no bibliography and no bibliography entry.  No retained line contains a figure environment, an image-inclusion command or drawing code, so no image is delivered and the package declares no asset dependency.  Editorial formatting only, in addition: an AMS \texttt{amsart} preamble that restates the article's own declarations and macros used by the retained lines, namely the environments \texttt{claim} on separate counters, together with the standard packages those lines need.  The retained environments print in this document as Claim 1 in order of appearance, whereas the article prints the same items with the numbering of its own full document.  The complete native source of the article as distributed in the arXiv e-print for this exact version is bundled unchanged as \texttt{sources/200454/source0.tex}.
\begin{claim}
    There is a balanced component of $A$ or $B$ having precisely $2$ vertices in $S$.
\end{claim} 

\begin{proof}
Suppose that the claim is false.
We use $r_G$ to denote the rank function of $M(G)$ ($r$ is used for the rank function of $M$). We define $d$ to be $(r(A) - r_G(A)) + (r(B) - r_G(B))$. Then the equation for the order of $(A,B)$ tells us that
\[
1 = r(A) + r(B) - r(E(G)) = (r_G(A) + r_G(B) - r_G(E(G))) - 1 + d \ge d,
\]
since $G$ is $2$-connected. So one of $r(A) - r_G(A)$ or $r(B) - r_G(B)$, without loss of generality the second, must be zero. Thus $B$ is balanced. Since every component of $A$ meets $S$ in at least $2$ vertices we have $c_A \le |S|/2$. 
If some component of $B$ meets $S$ in $2$ vertices, then the statement holds, so we may assume that every component of $B$ meets $S$ in at least $3$ vertices,  $c_B \le |S|/3$. Using the above formula for the order of $(A,B)$ once more, we have
\begin{align*}
1 &= (r_G(A) + r_G(B) - r_G(E(G))) - 1 + d \\
&= |V(A)| - c_A + |V(B)| - c_B - |V(G)| + d \\ 
&= |S| - c_A - c_B + d \\
&\ge |S| - |S|/2 - |S|/3 + d \\
&= |S|/6 + d \\
&> d.
\end{align*}
So $d = 0$ and thus $A$ is also balanced and by the above claim we have $c_A \le |S|/3$. Substituting this into the above calculation we have
\[
1 = |S| - c_A - c_B \ge |S| - |S|/3 - |S|/3 = |S|/3,
\]
so $S$ contains at most $3$ vertices. Since both $A$ and $B$ are balanced, every cycle in $\mathcal{F}$ must meet both $A$ and $B$ and so must contain at least $2$ of these $3$ vertices, contradicting the assumption that $\mathcal{F}$ is non-degenerate.
\end{proof}

\end{document}
